\begin{table}[ht]
\centering
\caption{Average rally length reached in each window, the absolute gain between them, and the relative improvement. Relative improvement normalises by each session's own first-window mean, so it rewards a low starting point.}
\label{tab:level}
\begin{tabular}{lrrrr}
\toprule
Group & Window 1 & Window 2 & Absolute gain & Relative improvement (\%) \\
\midrule
CFL-4 & 1.325 & 2.482 & +1.157 & +106.2 \\
CFL-3 & 0.988 & 1.439 & +0.450 & +75.6 \\
DP-5 & 0.885 & 1.189 & +0.304 & +70.9 \\
DP-10 & 0.825 & 1.179 & +0.354 & +66.9 \\
CFL-1 & 0.873 & 1.019 & +0.147 & +47.2 \\
CFL-2 & 0.835 & 1.011 & +0.176 & +40.7 \\
DP-20 & 0.835 & 0.980 & +0.145 & +40.0 \\
DP-2 & 0.852 & 0.957 & +0.105 & +49.0 \\
HCC & 0.637 & 0.880 & +0.243 & +57.1 \\
MCC & 0.700 & 0.847 & +0.147 & +39.3 \\
AIF-1 & 0.817 & 0.801 & -0.016 & +15.8 \\
\bottomrule
\end{tabular}
\\[4pt]\footnotesize{AIF-1 illustrates the limitation of the relative measure: its absolute change is negative ($-0.016$) while its mean relative improvement is positive ($+15.8\%$), because the mean of per-session ratios is not the ratio of means and is skewed by sessions with small first-window denominators.}
\end{table}
